\section{Introduction}
\label{sec:intro}

The integration of spectroscopic imaging with deep learning has produced
a class of compositional models that has become the de facto standard
across infrared tissue classification, hyperspectral remote sensing,
Raman imaging, and mass spectrometry imaging. A canonical pipeline
processes a hyperspectral input $X \in \R^{H \times W \times S}$ in two
stages: a per-pixel spectral reduction $f_\thetap: \R^S \to \R^K$
followed by a spatial model $g_\phip: \R^{K \times H \times W} \to
\R^{N_{\mathrm{cls}} \times H \times W}$. End-to-end joint training of
this pipeline is overwhelmingly the default — adopted by foundational
work in FTIR histology~\cite{berisha2019deep}, recent state-of-the-art
spatial-spectral ViTs for prostate tissue
classification~\cite{oleary2026spatial}, and essentially all
hyperspectral remote sensing systems.

\paragraph{An empirical paradox.}
Despite the prevalence of this paradigm, an unexpected and persistent
pattern has emerged across multiple subfields. Frozen pretrained
spectral encoders dramatically outperform end-to-end joint training
on held-out test data. In our motivating tissue classification setting,
a learned-from-scratch linear $f_\thetap$ trained end-to-end achieves
$0.675$ test F1; freezing a pretrained $f_\thetap$ yields $0.84$--$0.95$
test F1. Most strikingly, replacing $f_\thetap$ with \emph{frozen
random weights} yields $0.70$ test F1 --- outperforming the joint
trained model. Recent independent observations
\cite{oleary2026spatial,berisha2019deep} have noted, in different
terms, that aggressive spectral compression to $16$ features has
negligible impact on spatial-spectral neural network performance, and
have concluded that ``tissue classification only needs a small set of
spectral features.''

We argue the opposite. The reason end-to-end joint training fails is
not that spectral information is unimportant. It is that
\emph{end-to-end optimization fails to extract that information}, even
when the spectral signal is rich and discriminative. The problem is in
the training procedure, not the data.

\paragraph{This paper.}
We provide a rigorous explanation for the failure of joint
spectral-spatial training, grounded in the geometry of the loss
landscape and the dynamics of stochastic gradient descent on
compositional architectures with asymmetric capacity. Two main
theorems form our contribution.

\begin{itemize}
    \item \textbf{Theorem~\ref{thm:hessian} (Capacity-Induced Module
          Curvature Disparity).} We prove that the ratio of the
          spatial to the spectral Gauss-Newton block's top eigenvalue
          grows linearly with the spatial module's width $M$
          (equivalently, as $\sqrt{\Cg/\Cf}$ for architecture
          families with $\Cg = \Theta(M^2)$): as the spatial module
          widens, curvature concentrates in it while the spectral
          block's curvature stays capped by a width-independent
          constant. The proof is elementary: an explicit witness
          direction in the classifier-head weights forces the
          width-linear growth, and an operator-norm bound caps the
          spectral block; the rate matches the mean-field law of
          Karakida et al.~\cite{karakida2019universal} in the model
          class where the latter applies.
    \item \textbf{Theorem~\ref{thm:twoscale} (Finite-Time Spectral
          Starvation).} Under an explicit shortcut-fitting hypothesis
          compatible with cross-entropy dynamics, the spectral
          module's displacement from initialization by the time the
          training residual reaches level $\delta$ is at most
          $\mathcal{O}(\varepsilon \log(1/\delta))$, where
          $\varepsilon$ is the ratio of the (capped) spectral gradient
          scale to the shortcut fitting rate. As spatial capacity
          grows, the shortcut fits faster while the spectral gradient
          stays capped, so the spectral module is pinned near
          initialization for the entire useful duration of training.
          The formulation is finite-time by design: for unregularized
          cross-entropy on separable data no finite stationary point
          exists~\cite{soudry2018implicit}, and the starvation
          phenomenon is a training-time effect, not an asymptotic
          one.
\end{itemize}

\paragraph{A real-time diagnostic.}
The pathology described by Theorems~\ref{thm:hessian} and
\ref{thm:twoscale} is invisible to standard loss curves --- training
and validation loss both decay smoothly while the spectral module is
silently starved. To make the pathology observable in practice we
introduce the \emph{Effective Gradient Ratio} (EGR), a scalar
diagnostic computable inside any training loop:
\[
    \EGR(t) = \frac{\norm{\nabla_\thetap \loss(\thetap_t, \phip_t)}}
                   {\norm{\nabla_\phip \loss(\thetap_t, \phip_t)}}.
\]
Proposition~\ref{prop:egr_monotonic} establishes that the spectral
gradient norm---the EGR's numerator---decays polynomially while the
loss is being fitted. Empirically, the depth of the EGR collapse
scales monotonically with capacity ratio, providing a quantitative,
training-time observable for gradient starvation.

\paragraph{Prescription.}
The theory yields a direct prescription: freeze the spectral module,
or inject a fixed residual baseline. Freezing eliminates the timescale
gap by removing $\thetap$ from the optimization entirely. A residual
baseline ($Z = f_\thetap(X) + h(X)$ for a fixed $h$, e.g.\ PCA) is a
softer intervention that breaks the shortcut without precluding spectral
fine-tuning. We discuss both in Section~\ref{sec:discussion}.

\paragraph{Universality.}
Although our motivating application is infrared tissue classification,
the mechanism we identify applies to any compositional architecture
with capacity asymmetry between sequential modules. We verify in
Section~\ref{sec:experiments} that the same pattern manifests across
linear, CNN, and Vision Transformer spatial models on synthetic
hyperspectral problems, indicating the phenomenon is architecture-
agnostic. We expect the same mechanism to govern frozen-encoder
practices in Raman spectroscopy, mass spectrometry imaging, NIR
analysis, and beyond.

\paragraph{Relation to prior work.}
The closest prior work is in multimodal learning, where Wang et
al.~\cite{wang2020what} and Peng et al.~\cite{peng2022balanced}
identified ``modality competition'' and proposed gradient blending and
modulation interventions. Our setting is unimodal (a single
hyperspectral input) but the mechanism is structurally analogous: the
spatial and spectral subspaces of the input act as competing
``modalities,'' with the spatial module monopolizing the gradient
budget. We discuss this connection further in
Section~\ref{sec:related}.

\paragraph{Outline.}
Section~\ref{sec:related} reviews related work. Section~\ref{sec:setup}
establishes notation. Section~\ref{sec:thm1} proves
Theorem~\ref{thm:hessian}. Section~\ref{sec:thm2} proves
Theorem~\ref{thm:twoscale}. Section~\ref{sec:egr} formalizes the EGR
diagnostic. Section~\ref{sec:experiments} validates both theorems and
the diagnostic on synthetic spectral-spatial problems with linear, CNN,
and ViT spatial models. Section~\ref{sec:real} briefly connects the
synthetic findings to real FTIR/QCL tissue classification data.
Section~\ref{sec:discussion} discusses prescriptions, limitations, and
open problems.
